\clearpage

\lehead[]{\normalfont\sffamily\hspace*{-2.00cm}\textcolor{white}{\colorbox{lightblue}{\parbox[c][0.70cm][b]{1.60cm}{
\makebox[1.60cm][r]{\thechapter}\\ \makebox[1.60cm][r]{ÜBUNG}}}}\hspace{0.17cm}\textcolor{lightblue}{\chaptertitle}}
\rohead[]{\textcolor{lightblue}{\chaptertitle}\normalfont\sffamily\hspace*{0.17cm}\textcolor{white}{\colorbox{lightblue}{\parbox[c][0.70cm][b]{1.60cm}{\thechapter\\
ÜBUNG}}}\hspace{-2.00cm}}
%\chead[]{}
\rehead[]{\textcolor{lightblue}{AvHG, Inf, My}}
\lohead[]{\textcolor{lightblue}{AvHG, Inf, My}}

\section{Threads -- Übungen}

\subsection{Aufgabe 1: Primzahl-Suche ohne Einfrieren der Oberfläche}

Wenn ein Programm eine langwierige Berechnung im Haupt-Thread (dem Event
Dispatch Thread von Swing) ausführt, friert die gesamte Benutzeroberfläche ein,
bis die Berechnung beendet ist.

\begin{compactenum}[a)]
\item Erstelle ein \myClass{JFrame} mit zwei Buttons (\myUserInput{Start} und
\myUserInput{Abbrechen}) sowie einem \myClass{JLabel} zur Statusanzeige.

Schreibe eine Methode, die zählt, wie viele Primzahlen es zwischen $2$ und
$200.000$ gibt (prüfe dazu für jede Zahl $z$ in einer Schleife, ob sie durch
irgendeine Zahl von $2$ bis $z - 1$ ohne Rest teilbar ist -- absichtlich ohne
mathematische Optimierung, damit die Berechnung einige Sekunden dauert!).

Rufe diese Methode zunächst direkt im \myClass{ActionListener} des
\myUserInput{Start}-Buttons auf. Beobachtung: Während das Programm rechnet,
lässt sich weder das Fenster verschieben noch der \myUserInput{Abbrechen}-Button
anklicken!

\item Lagere die Primzahl-Suche nun in eine eigene Klasse aus, die von
\myClass{Thread} erbt. Während der Thread rechnet, soll das \myClass{JLabel}
regelmäßig den aktuellen Zwischenstand anzeigen (z.\,B.\ alle $5.000$ geprüften
Zahlen).

\item Ergänze die Thread-Klasse um eine boolesche Variable
\lstinline|abgebrochen|, sodass der Benutzer die laufende Berechnung jederzeit
über den \myUserInput{Abbrechen}-Button sauberen beenden kann.
\end{compactenum}


\subsection{Aufgabe 2: Wettrennen der Threads}

Es soll ein Wettrennen zwischen vier Läufern (oder vier Fahrzeugen)
programmiert werden, die jeweils durch einen eigenen \myClass{Thread}
gesteuert werden.

\begin{compactenum}[a)]
\item Erstelle ein \myClass{JFrame} mit einem \myUserInput{Start}-Button, einem
\myUserInput{Reset}-Button und einer Zeichenfläche (\myClass{JPanel}), auf der
vier Bahnen mit je einem Läufer (z.\,B.\ als farbiger Kreis oder Rechteck)
sowie eine Ziellinie bei $x = 500$ gezeichnet werden.

\item Programmiere eine Klasse \myClass{LaeuferThread}, die von
\myClass{Thread} erbt. Jeder Läufer bewegt sich in einer Schleife jeweils um
einen zufälligen Schritt (z.\,B.\ $1$ bis $5$ Pixel) nach rechts, stößt ein
\lstinline|repaint()| an und legt sich danach für eine zufällige Zeitspanne
(z.\,B.\ zwischen $10$ und $40$ Millisekunden) schlafen.

\item Beim Klick auf \myUserInput{Start} werden alle vier Threads gleichzeitig
gestartet. Sobald der erste Läufer die Ziellinie erreicht oder überschreitet,
wird der Sieger im Fenster angezeigt, und alle anderen Läufer-Threads sollen
sofort stehen bleiben (Tipp: Verwende eine gemeinsame
\lstinline|static volatile boolean|-Variable \lstinline|rennenBeendet|, die
alle Threads in ihrer Schleifenbedingung abfragen).
\end{compactenum}


\subsection{Aufgabe 3: Ampelschaltung mit Nachtmodus}

Bei einer Ampel dauern die verschiedenen Phasen unterschiedlich lange (z.\,B.\
Rot: $4\,\text{s}$, Rot-Gelb: $1\,\text{s}$, Grün: $4\,\text{s}$, Gelb:
$1{,}5\,\text{s}$). Mit einem Thread lässt sich eine solche Abfolge dank
\lstinline|Thread.sleep()| direkt von oben nach unten aufschreiben.

\begin{compactenum}[a)]
\item Programmiere ein \myClass{JFrame} mit einer gezeichneten Ampel (drei
Kreise für Rot, Gelb und Grün) sowie zwei Buttons: \myUserInput{Tagbetrieb} und
\myUserInput{Nachtmodus} (optional zusätzlich \myUserInput{Aus}).

\item Programmiere eine Klasse \myClass{AmpelThread}, die von \myClass{Thread}
erbt. Im Tagbetrieb durchläuft der Thread endlos die vier Ampelphasen mit den
oben genannten Wartezeiten und sorgt nach jedem Phasenwechsel für ein
Neuzeichnen der Ampel.

\item Wenn der Benutzer auf \myUserInput{Nachtmodus} klickt, soll der laufende
Tag-Thread sauber beendet und stattdessen ein Blink-Thread gestartet werden,
der das gelbe Licht im Takt von $500$ Millisekunden ein- und ausschaltet.
\end{compactenum}

